% Integration notes: requires amsmath, amssymb, amsthm and theorem/proposition.
% Approximately 3--5 article pages after house-style compression if needed.

\section{From terminal certification to global complexity}
\label{sec:global-bridge}

The target considered here is
\[
 n^{O(\log n)}=\exp\bigl(O((\log n)^2)\bigr),
\]
where $n\geq2$ is the original diagram size.  The broader term
\emph{quasi-polynomial} also permits higher fixed powers of $\log n$ in
the exponential.  We distinguish these bounds explicitly.  None of the
global geometric hypotheses below has been proved for every native
hierarchy state in this package.

The new terminal algorithm removes one local obstruction: a source cut
with disk parallelity bases can be accepted through a signed group proof
even when its unsigned stopping core is large.  It does not produce the
surfaces leading to that cut, control the number of candidates, or
transport a boundary pattern through subsequent cuts.  The following
results identify precise ways to control those remaining costs.

\subsection{Amortized branching, including long repair sequences}

Consider the full finite computation tree of a hierarchy search.  A
node can be a genuine choice of cutting surface, a conjunction requiring
several cut pieces to be solved, or a deterministic repair.  Count every
recursive subproblem, including unsuccessful candidates.  Cheap rejected
candidates may instead be charged to the node's local work, provided
their entire generation and rejection cost is bounded there.

\begin{proposition}[An amortized branching bound]
\label{prop:entropy-search}
Suppose the following conditions hold for inputs of size $n$.
\begin{enumerate}
\item The search has a proved sound and complete hierarchy specification.
      Source creation, candidate generation, admissibility checking,
      marked transport and terminal discovery and verification cost at most $n^C$
      per node and per emitted recursive edge.  All required encoded
      states and local source proofs have at most $n^C$ bits.
\item Every root-to-leaf path has at most $n^C$ nodes.
\item If $b(v)\geq1$ is the number of children of a nonleaf node, then
      every root-to-leaf path $P$ satisfies
      \[
         \sum_{v\in P\text{ nonleaf}}\log_2 b(v)
                      \leq A(\log_2 n)^2.
      \]
\end{enumerate}
Then the total search, including unsuccessful candidates, costs
$n^{O(\log n)}$ bit operations.

A local sufficient condition for the third hypothesis is a certified
integer mass $1\leq M(v)\leq n^D$, nonincreasing along deterministic
edges, for which every branching edge $v\to w$ satisfies
\begin{equation}
 \log_2 b(v)\leq
 K\bigl((\log_2 M(v))^2-(\log_2 M(w))^2\bigr)+e(v,w),
 \label{eq:quadratic-log-charge}
\end{equation}
where all charges are nonnegative and their sum on each path is
$O((\log n)^2)$.
\end{proposition}
\begin{proof}
Assign equal conditional probability $1/b(v)$ to every edge leaving
each nonleaf node.  A leaf $\ell$ has probability
\[
 p_\ell=\prod_{v\in P_\ell}b(v)^{-1}\geq
             2^{-A(\log_2 n)^2}.
\]
Leaf probabilities sum to one.  Hence there are at most
$2^{A(\log_2 n)^2}$ leaves.  Every node lies on a root-to-leaf path, so
the node count is at most the maximum path length times the number of
leaves.  The edge count is smaller than the node count.  Multiplying by
the stated local bit cost proves the bound.

For the sufficient condition, sum~\eqref{eq:quadratic-log-charge} along
a path.  The squared-log terms telescope; deterministic edges have zero
branching cost and nonnegative potential drop.  The initial potential
is at most $D^2(\log_2 n)^2$, and the accumulated charges have the stated
bound.  This gives the required path sum.
\end{proof}

For example, if every branching child has mass at most half its parent
and $b(v)\leq M(v)^K$, then~\eqref{eq:quadratic-log-charge} holds with
zero extra charge.  Indeed, with $x=\log_2 M(v)\geq1$, the potential
drop is at least $x^2-(x-1)^2=2x-1\geq x$.  Thus a balanced
marked-cut theorem with polynomially many local alternatives is enough,
even if polynomially long deterministic repairs occur between its
branching steps.  The mass need not be the number of represented normal
disks; it must be a source-certified quantity with the proved decrease.

The path-length and encoding hypotheses have independent roles.  A
branching estimate does not bound an arbitrarily long unary computation,
and a bounded number of handle families does not bound the bit length of
their endpoints or the number of newly split marking intervals.  The
terminal cost is polynomial when the current $t$, profile list, integer
bit mass and source-proof size are polynomial in $n$.  This fact does
not establish the corresponding transport bounds for a successor.

\subsection{An exact small terminal key supplied by the new compiler}

The signed disk-base compiler has a useful compression property beyond
its one-shot runtime.  Here $W$ is the authenticated compact orientable
PL whole cut, every component has nonempty boundary, and every parallelity
base is a disk.  Write $a,m$ for its initial generator and relator counts
and $h_3$ for its source three-handle count.  The source ledger is
\[
 \chi(W)=c-a+m-h_3,
 \qquad r_0:=c-\chi(W)=a-m+h_3.
\]
The whole-cut criterion says exactly that $W$ is a union of handlebodies
when its component-group free product is free of rank $r_0$.

\begin{proposition}[Exact residual terminal kernel]
\label{prop:terminal-key}
Apply any independently verified finite Tietze prefix to the compiled
presentation.  Suppose $s$ surviving generators occur in no remaining
relator.  Remove these generators from the displayed presentation and
call the remaining presentation $Q$.  Put $r_*=r_0-s$.  Then
\begin{equation}
 W\text{ is a union of handlebodies}\quad\Longleftrightarrow\quad
 G(Q)\cong F_{r_*},
 \label{eq:residual-terminal}
\end{equation}
where the right side is false for $r_*<0$.

Let $k\geq1$ bound the number of generators, relators and literal letters
of $Q$ in total.  Impossible target ranks can be rejected immediately.
For $0\leq r_*\leq k$, compact relabelling gives a key $(Q,r_*)$ with
$O(k\log(k+2))$ bits, and there are
$2^{O(k\log(k+2))}$ such keys.  Normal-coordinate magnitudes and
isolated-ball multiplicities do not occur in this terminal key.
\end{proposition}
\begin{proof}
The full current group is $G(Q)*F_s$.  If $G(Q)\cong F_{r_*}$, its
free product with $F_s$ has rank $r_0$, giving the forward terminal
conclusion.  Conversely, if $G(Q)*F_s\cong F_{r_0}$, then $G(Q)$ is an
embedded finitely generated subgroup and hence free.  Abelianization
shows that its rank is exactly $r_0-s$.  This proves the equivalence.

If the required rank is negative the test is already negative.  A free
group presented on at most $k$ generators cannot have rank exceeding
$k$, so values outside $0\leq r_*\leq k$ can be separated immediately.
Encode the words with row separators over a signed alphabet of at most
$k$ generators.  The encoded string has $O(k)$ symbols from an alphabet
of $O(k)$ symbols, together with the generator count and target rank.
This proves both the bit bound and the counting bound.  The source
compiler's Euler formula shows that isolated balls cancel between $c$
and $\chi(W)$, while a verified Tietze prefix eliminates all dependence
on the original geometric labels for this algebraic predicate.
\end{proof}

This is an exact reduction of the terminal predicate, not an algorithm
for deciding arbitrary group freeness.  The present signed policy is
still one-sided.  Equal syntactic keys permit reuse of an algebraic
certificate after replaying the source compiler and the relabelling map.
Each source occurrence still needs its own authenticated binding.  A
cached greedy refusal remains a refusal of that policy.

The initial active key has $O(t\log(t+2))$ bits.  If future normalizations
guarantee residual size
\[
       k\log(k+2)=O((\log n)^2),
\]
there are at most $n^{O(\log n)}$ distinct relabelled syntactic terminal
keys.  In particular $k=O(\log n)$ gives the smaller bound
$n^{O(\log\log n)}$.  The parameter is the remaining signed
presentation size, not the old unsigned deferral number: the earlier
$q$-block geometric stopping family reduces to the empty residual even
though its raw deferral is $q$.

This cache cannot identify hierarchy states.  Equal abstract groups and
Euler data do not identify peripheral curves, marked faces, or admissible
next surfaces.  Nor does a small number of terminal keys bound the cost
of generating the geometric states that reach them.

\subsection{A precise marked-interface route to memoization}

There is a distinct route when many search histories give the same
future marked problem.  The required equivalence is stronger than
isomorphism of the terminal groups.

\begin{proposition}[Source-replayable interface quotient]
\label{prop:marked-interface-quotient}
Suppose a sound and complete marked hierarchy specification admits
keys with the following properties, where $1\leq k\leq n^C$.
\begin{enumerate}
\item A key has at most $k$ active atoms compactly labelled in
      $\{1,\ldots,k\}$, at most $O(k)$ incidence
      and cyclic-order records, and at most $O(k)$ references to a fixed
      source dictionary of at most $n^C$ exact objects.  The dictionary
      has polynomial total bit mass.  Each record has bounded arity and
      uses only atom labels, dictionary labels and a fixed finite alphabet.
\item Equal keys have source-replayable equivalences preserving the
      entire continuation problem: boundary pattern, orientations,
      peripheral data, admissibility and the correctness of all cuts.
      A representative key has at most $n^C$ candidate rules with at most
      $n^C$ child keys each, generated and checked in polynomial time.
      These rules are complete for every source realization of that key.
\item Terminal acceptance and its certificate are computed in polynomial
      time per key, or the terminal certificates are included among the
      polynomially generated candidate rules.  Merely knowing that a
      short certificate exists does not satisfy this hypothesis.
      Each terminal proof and each local source equivalence has a
      polynomial-size, polynomially checkable representation.  The
      source-lifting verifier can compose these objects as a proof DAG
      in time polynomial in $n$ per DAG node.  Source reconstruction,
      coordinate arithmetic, reindexing and marked transport are included
      in these bounds; the proof does not require exponentially expanding
      shared histories.
\end{enumerate}
Then the hierarchy can be solved in
\[
 \exp\bigl(O(k(\log n+\log(k+2)))\bigr)\operatorname{poly}(n)
\]
time.  In particular $k=O(\log n)$ yields $n^{O(\log n)}$.
\end{proposition}
\begin{proof}
There are at most
\[
 (n^C)^{O(k)}(k+1)^{O(k)}
   =\exp\bigl(O(k(\log n+\log(k+2)))\bigr)
\]
keys.  Generate the reachable quotient once, reusing every previously
seen key.  The number of rules and child references is the key count
times a polynomial in $n$.

A rule can require all its children to be accepted, while several rules
at one key represent alternative choices.  Thus acceptance is the least
fixed point of a finite monotone OR/AND rule system.  Maintain, for each
rule, the number of children not yet accepted; enqueue its parent when
that count reaches zero.  Each child reference is processed once.  Cycles
do not certify themselves: a key is accepted only after an accepting
terminal or a rule whose children were accepted earlier.  Back-pointers
therefore give an acyclic accepting proof, even if the explored quotient
contains cycles.  The completeness hypothesis identifies this fixed
point with the original finite-proof hierarchy problem.  The source
lifting hypothesis makes the resulting proof DAG valid for the original
marked input.  All generation, fixed-point and replay costs have the
stated bound.
\end{proof}

The common dictionary is a substantive requirement.  Its entries may be
very large integers, provided there are only polynomially many exact
entries and the total bit mass is polynomial.  If each of $k$ marked
coordinates instead varies freely over $B$ bits, the key bound becomes
$2^{O(k(B+\log k))}$.  With $k=\Theta(\log n)$ and
$B=\Theta((\log n)^2)$, this gives $\exp(O((\log n)^3))$, not the
target $n^{O(\log n)}$.  Proving that successor endpoints stay in a
small dictionary, or proving an equally strong exact normal form, is
part of the geometric work.

\subsection{The concrete next bridge and the exponent traps}

The most direct continuation from the implemented operator is a
\emph{marked balanced-cut theorem}: after source-authenticated
parallelity compression, either the entire piece has disk bases and a
bounded signed terminal certificate, or a polynomially generated,
complete family of marked cuts satisfies the mass charge in
\eqref{eq:quadratic-log-charge}.  Every chosen successor must retain
polynomial bit size and an exact transport of its inherited pattern.
The new finite group compiler enlarges the acceptable terminal class;
the missing theorem is the guaranteed arrival in that class or the
controlled progress of every recursive alternative.  A proof that a
particular word policy always succeeds on the guaranteed terminal class
would be required; its current Fibonacci counterexample prevents a
general freeness-completeness claim.

There are three common accounting failures.  Polynomial local work at
polynomial branching depth can still cost $2^{n^{\Omega(1)}}$ even with
two children per level.  A branching factor $n^{O(\log n)}$ at each of
$O(\log n)$ levels gives
$n^{O((\log n)^2)}=\exp(O((\log n)^3))$.  By contrast, multiplying a
total call count $n^{O(\log n)}$ by a per-call cost
$n^{O(\log n)}$ preserves $n^{O(\log n)}$: exponents add for total
work, whereas branch choices multiply along a path.  Finally, counting
only successfully emitted surfaces omits the potentially dominant cost
of unsuccessful surface searches.  Both propositions above explicitly
charge those searches and require a complete candidate mechanism.

\paragraph{Recheck of the native constants.}
The bounds $a\leq74t$ and $m,L\leq78t$ remain valid for the present
source format.  Their word-length argument uses the source validator's
exclusion of reverse-edge identifications and its circle-link condition:
an interior edge-link traversal visits each local tetrahedral edge
occurrence exactly once, so all wedge lists total at most $6t$; at most
two endpoint ranks double this to $12t$.  The additional annular words
use at most the $66t$ authenticated ports once in total.  These bounds
do not constrain the full base-profile list, integer bit mass, source
proof length, transported marking intervals, or the tetrahedron count
of later successor states.  Those quantities cannot be absorbed into
the constants $74$ and $78$.
